\section{实现}
\label{sec:implementation}

\subsection{架构与可复现性}
\label{sec:architecture}

原型用 Python 3.11 实现，依赖锁定在 \texttt{requirements.txt}，随机种子与配置文件
一并提交。核心分流逻辑只有一个小模块；审阅界面是命令行工具，便于在无图形环境的服务器
上复现。图~\ref{fig:pipeline} 中的每个阶段都有独立测试，审计记录写成 JSON Lines，
一行一个动作。

下面的函数是分流逻辑的最小示例：模型分数高于阈值且存在证据片段时自动通过，否则进入
人工审阅。

\inputminted[fontsize=\footnotesize]{python}{code/example.py}

\subsection{组件与职责}
\label{sec:components}

表~\ref{tab:components} 列出主要组件及其职责。组件之间的接口是显式的：分类器只返回
标签与分数，声明抽取器只返回原文片段与位置，审计记录只追加、不修改。

\begin{table}[htbp]
  \centering
  \formaltable
  \begin{tabularx}{\textwidth}{|B|Y|}
    \hline
    \textbf{组件} & \textbf{职责} \\ \hline
    文档接入 & 读取文档与元数据，检查格式与授权字段。 \\ \hline
    分类器 & 给出主题标签与置信度；不直接拒绝文档。 \\ \hline
    声明抽取器 & 返回支持标签的原文片段与位置。 \\ \hline
    人工审阅 & 接受、修改或拒绝提议；最终标签由人决定。 \\ \hline
    审计记录 & 追加模型输出、审阅动作、时间与理由。 \\ \hline
  \end{tabularx}
  \caption{系统组件与职责。}
  \label{tab:components}
\end{table}

